\section{Manufacturing an eight-summary finite-field terminal}\label{sec:eightmoments}
The remaining field-coordinate scan can be removed from the complexity
theorem on a further acquired family. The number of module coordinates
is still unbounded. The extra certificate constructs a residue normal form
\begin{equation}\label{eq:pairnormal}
 H_0=\gamma\sum_{i=1}^r
 \left(X_i^3+X_iY_i^2+Y_i^3+
 Z(X_iY_i+a_iX_i^2+b_iY_i^2)\right)+\tau Z^3,
\end{equation}
with $\gamma\ne0$, $a_i,b_i\in\F_2$, and arbitrary $\tau\in F$.
The basis, scalar and pair types are manufactured. This includes $J_r$,
arbitrary invertible field coordinate changes and phase scaling, then
arbitrary higher nilpotent cubic coefficients, compatible quotient noise,
and the outer raw binary affine disguise. Require $d\ge3$ and $3\nmid d$.

\begin{theorem}[All-depth acquired finite-moment counter]\label{thm:eightmoments}
The preceding family has deterministic polynomial-time exact counting,
uniformly in $b=2r+1$, $d$, and $e$, with $n=bed$ original input bits.
A conservative bit-time bound is
\[
 O(Ln^2+n^{12})+\widetilde O(d^{31}),
\]
and working space is $O(L+n^6)+\widetilde O(d^{20})$.
The tilde terms occur only at a depth-one terminal. That terminal uses
seven explicit cubic point counts in at most eight field variables and
one binary Gaussian, followed by exact linear reconstruction.
\end{theorem}
The generic fixed-variable point-count engine is the established
Lauder--Wan algorithm \cite{lauderwan}; it is not implemented in this
source distribution. Acquisition, coefficient verification, the eight
reductions and reconstruction are implemented. The executed reference
backend explicitly scans a bounded field and is distinguished from the
polynomial-time theorem.

\subsection{Acquiring binary structure by semilinear descent}
Acquire the contact chart as before. Its affine gradient has a unique
binary solution for this family; solve it and change $w\mapsto w+az$.
The $z^2L_1$ term disappears. Write the centered polynomial as
$G_3(w)+zQ_2(w)+\tau z^3$, with square matrix $M$ and polar matrix
$\Omega$. Set $\sigma(w)=w^{[2]}$ and form
\[
 S=\Omega^{-1}M\sigma,\qquad K=S^d.
\]
The latter is field-linear and is computed by $d$ semilinear compositions.
Check $K^2+K+I=0$. Let $a\in\{1,2\}$ satisfy $ad\equiv1\pmod3$ and
put $D=K^a$, $F_*=D^{-1}S$. Acquire
$U=\ker_{\F_2}(F_*-I)$. Check that $\dim_{\F_2}U=b-1$ and that its
binary basis is a field basis of the transverse space.

In that basis, check that every entry of $\Omega$ is zero or one common
nonzero field scalar $\gamma$. The normalized form is binary symplectic.
Compute $D$ on $U$, check that its entries are binary and that it
preserves this form. These are polynomial-size identity tests, with no
enumeration of field elements or possible bases.

Manufacture $D$-invariant symplectic planes. Since $D^2+D+I=0$, the
binary quadratic $v\mapsto\Omega(v,Dv)/\gamma$ has polar form
$\Omega/\gamma$. It is nonzero, so a vector with value one is found
among a basis and its pairwise sums. The plane spanned by $v,Dv$ is
nonsingular, and its orthogonal complement is $D$-invariant. Repeat.
Every acquired plane has blocks
\[
 \Omega/\gamma=\begin{pmatrix}0&1\\1&0\end{pmatrix},\qquad
 D=\begin{pmatrix}0&1\\1&1\end{pmatrix}.
\]
Because the basis lies in $U$, $S=D\sigma$ there, hence $M=\Omega D$.
This gives exactly the repeated cubic blocks in \eqref{eq:pairnormal}.
Read the square coefficients of $Q_2$ and require their quotients by
$\gamma$ to be binary. Finally verify the \emph{entire ordinary
polynomial identity}, not a sample of Boolean values.

To prove completeness under hidden coordinates, start in the displayed
family. After acquiring the contact line and transverse center, the old
and new coordinates are $w_{\rm old}=Rw$, $z_{\rm old}=tz$, with $R$
invertible and $t\ne0$. If $D_0$ is the sum of the displayed order-three
blocks and the original scalar is $\gamma_0$, then
\[
 S=t^{-1}R^{-1}D_0\sigma R,\quad
 K=R^{-1}D_0^dR,\quad D=R^{-1}D_0R,\quad
 F_*=t^{-1}R^{-1}\sigma R.
\]
The scalar norm of $t$ is one. Therefore
$U=R^{-1}t\F_2^{b-1}$ and the common symplectic scale is
$\gamma_0t^3$. Every test succeeds. The acquired planes may differ
from the hidden ones, but the normalized square coefficients remain
binary and the full coefficient identity remains valid. The earlier
centroid and quotient-fit proof handles the outer binary affine disguise.

The rotation $(X,Y)\mapsto(Y,X+Y)$ preserves the cubic block and sends
square type $(a,b)$ to $(b,1+a+b)$. It cycles $(0,0),(0,1),(1,0)$ and
fixes $(1,1)$. Thus only two counting types remain, with acquired
multiplicities $r_+$ and $r_-$.

\subsection{Symmetry gives eight joint Gaussian states}
For fixed $z\in F$, define
\[
 A_\pm(z)=\sum_{x,y\in F}(-1)^{\Tr\{\gamma[
 x^3+xy^2+y^3+z(xy+\delta_\pm(x^2+y^2))]\}},
 \quad\delta_+=0,\quad\delta_-=1.
\]
The phases are binary quadratic on $2d$ bits with the same polar matrix.
Squaring their common polar-kernel equations gives
\begin{align*}
 \gamma x^4+x+\gamma y^4+\gamma z^2y^2&=0,\\
 x+\gamma y^4+y+\gamma z^2x^2&=0.
\end{align*}
The ordinary Jacobian is
$\left(\begin{smallmatrix}1&0\\1&1\end{smallmatrix}\right)$.
All geometric solutions are isolated and simple; the degree bound gives
at most sixteen. The common binary polar nullity $k$ is therefore one
of $0,2,4$. Each Gaussian is zero or $\pm q2^{k/2}$.

The minus phase is invariant under $C(x,y)=(y,x+y)$, with
$C^2+C+I=0$. Its radical is $C$-stable. The quadratic is linear there;
invariance and $v+Cv+C^2v=0$ force it to vanish on the radical.
On the nonsingular quotient, manufacture $C$-invariant symplectic planes.
In each plane spanned by $v,Cv$ with polar pairing one, invariance forces
the quadratic to equal one on all three nonzero vectors. Each plane has
negative Gaussian sign. Consequently
\[
 A_-(z)=(-1)^{d-k/2}q2^{k/2}\ne0.
\]
When $k=0$, $A_+$ is also nonzero. With $\varepsilon=(-1)^d$, the joint
state $(A_+,A_-)$ belongs to the eight-element set
\[
 \mathcal V=\{(\pm q,\varepsilon q),\ (0,-2\varepsilon q),
 (\pm2q,-2\varepsilon q),\ (0,4\varepsilon q),
 (\pm4q,4\varepsilon q)\}.
\]
Its size is independent of $b,d,e$; each integer entry has $O(d)$ bits.

\subsection{Eight manufactured summaries and exact reconstruction}
For $(u,v)\in\mathcal V$, define the signed multiplicity
\[
 C_{u,v}=\sum_{z:(A_+(z),A_-(z))=(u,v)}(-1)^{\Tr(\tau z^3)}.
\]
Use moments $M_{ij}=\sum_z(-1)^{\Tr(\tau z^3)}A_+(z)^iA_-(z)^j$
at the eight indices
\[
 (0,0),(0,1),(0,2),(1,0),(1,1),(1,2),(2,0),(2,1).
\]
Their evaluation matrix on $\mathcal V$ is invertible over the rationals.
Group the nodes by their three distinct second coordinates. The first
three moments recover the groups' signed masses; the next three recover
their first $u$-moments. In the first group $u^2=q^2$ identically.
Subtract its known contribution from the last two moments and solve the
two-by-two system for the remaining second $u$-moments. These determine
each zero-node mass and the separate masses at opposite nonzero values.
This proves uniqueness and supplies exact reconstruction. The determinant
magnitude of the eight-by-eight matrix is $2^{15}3^7q^{14}$.

The field count is then
\[
 Z_1=\sum_{(u,v)\in\mathcal V}C_{u,v}u^{r_+}v^{r_-},
\]
with $0^0=1$. Moment and reconstruction operands have $O(d)$ bits.
The final powers and counts have $O(bd)$ bits. No unit-cost operation
on an exponentially long integer is assumed.

Each positive-degree moment expands into at most three pair copies,
sharing $z$. It is the character sum of the explicit ordinary cubic
\[
 P_{ij}=\tau z^3+\gamma\sum_{\ell=1}^{i+j}
 [x_\ell^3+x_\ell y_\ell^2+y_\ell^3+
 z(x_\ell y_\ell+\delta_\ell(x_\ell^2+y_\ell^2))],
\]
where the first $i$ copies have $\delta_\ell=0$ and the last $j$ have
$\delta_\ell=1$. It has $2(i+j)+1\le7$ field variables. Manufacture
the Artin--Schreier equation $w^2+w+P_{ij}=0$, giving
\[
 M_{ij}=\#\{w^2+w+P_{ij}=0\}(F)-q^{2(i+j)+1}.
\]
These seven hypersurfaces have degree three and at most eight variables,
independently of the number of pairs. The zeroth moment is the binary
Gaussian $\sum_z(-1)^{\Tr(\tau z^3)}$.
Lauder--Wan's bound for fixed characteristic, degree and variable count
$s$ is $\widetilde O(d^{3s+7})$ time and
$\widetilde O(d^{2s+4})$ space \cite{lauderwan}. Taking $s\le8$ gives
the stated $d^{31}$ and $d^{20}$ bounds. The summaries are therefore
manufactured by fixed-dimensional computations, not assumed as advice.

The residue gradient in \eqref{eq:pairnormal} has the same zero binary
kernel as $J_r$ when $3\nmid d$. There is one projective critical ray
if $\tau=0$, and none if $\tau\ne0$. All ring stages are polynomial;
this terminal closes the remaining depth-one case and proves
Theorem~\ref{thm:eightmoments}.

\subsection{A sharp divisibility boundary in the same residue family}
The restriction $3\nmid d$ cannot simply be deleted while keeping arbitrary
higher nilpotent cubics. The very same $J_r$ residue already gives a
\#P-hard family at the fixed degree $d=3$ and depth $e=2$, with no quotient
quadratic. This is stronger than hardness in the broader contact class.

Let $p$ be a Boolean cubic on an even number $N\ge2$ of variables, and set
$r=N+1$, $b=2r+1=2N+3$. In $F=\F_8$, choose
$\alpha^3+\alpha+1=0$. A nonzero critical residue of $J_r$ has $Z\ne0$;
normalize $Z=1$. Its transverse equations give
\[
 X_i=Y_i^2,\qquad Y_i^4+Y_i^2+Y_i=0,
 \qquad Y_i=a_i\alpha+b_i\alpha^2,\quad a_i,b_i\in\F_2.
\]
The remaining derivative is $\sum_iX_iY_i=\sum_iY_i^3=0$.
The exact identity
\[
 Y_i^3=a_i\alpha+b_i\alpha^2+(a_i+b_i+a_ib_i)
\]
therefore makes the critical-ray constraints
\[
 \sum_i a_i=0,\qquad \sum_i b_i=0,\qquad \sum_i a_ib_i=0.
\]
The field-linear form $L_i=\alpha^2X_i+\alpha^4Y_i$ equals $a_i$ on
this chart. Manufacture the ordinary homogeneous cubic
\[
 G(X,Y,Z)=\sum_{S\in\operatorname{supp}(p)}
 Z^{3-|S|}\prod_{i\in S}L_i,
 \qquad H=J_r+tG\in F[t]/(t^2)[X,Y,Z].
\]
It has only polynomially many explicitly constructed coefficients.
On a critical chart point, $G=p(a_1,\ldots,a_N)\in\F_2$.

Because $r$ is odd, an even-parity vector $a$ cannot be the all-ones
vector. For $a\ne0$ the two equations on $b$ are independent and have
$2^{r-2}$ solutions; for $a=0$ there are $2^{r-1}$. The first $N$
coordinates determine the last $a$ uniquely. If
$Z_p=\sum_x(-1)^{p(x)}$ and $s_0=(-1)^{p(0)}$, then
\[
 N_{\rm crit}=2^{N-1}(2^N+1),\qquad
 W_{\rm crit}=2^{N-1}(Z_p+s_0).
\]
Summing all highest ring digits gives the factor $8^b$ and restricts the
residue to $\nabla J_r=0$. The origin contributes one. On each nonzero
ray, cubing permutes $F^\times$, so the radial sum is seven when $G=0$
and minus one when $G=1$, namely $3+4(-1)^G$. Consequently
\begin{equation}\label{eq:samepairhardness}
 Z(H)=8^{2N+3}
 \left[1+3\,2^{N-1}(2^N+1)
       +4\,2^{N-1}(Z_p+s_0)\right],
 \qquad n=12N+18.
\end{equation}
One exact count recovers $Z_p$ by polynomial-bit arithmetic. The full
binary centroid remains $F[t]/(t^2)$ by the earlier contact-centroid
proof, and arbitrary invertible binary affine disguises preserve the
count. Thus this is a raw acquired-algebra obstruction, not a supplied
decomposition counterexample.

For completeness, signed Boolean cubic counting is \#P-hard directly:
encode each bounded-fan-in Boolean gate by a quadratic equation
$g_j(x,y)=0$, include the accepted-output equation, and introduce a
fresh binary multiplier $u_j$. The phase $\sum_j u_jg_j$ is cubic.
Summing the multipliers enforces every equation and multiplies the
number of accepted computations by $2^{\#j}$. Gate values are uniquely
determined by the input. Add one unused bit if necessary to make $N$
even and divide its known factor out. This proves the asserted reduction
without assuming a lower bound on every large-kernel instance.

The executable manufactures two raw affine disguises. The 42-bit case
has ten critical rays, source signed count $-2$ and ring signed count
$14680064$; it independently scans all $2097152$ residue vectors.
The 66-bit case has 136 rays, source signed count zero and ring signed
count $3238405341184$; it independently scans 1024 chart candidates.
Both recover the source count through \eqref{eq:samepairhardness},
acquire a six-dimensional binary centroid, and pass 256 original-input
representation checks. Neither test enumerates its full ring domain.
An independent coefficient audit also covers the 90-bit construction.

\subsection{Execution boundary and prior machinery}
The source implements semilinear descent, plane acquisition, every
ordinary coefficient check, the eight-state reconstruction, and explicit
Artin--Schreier polynomial manufacture. Six disguised coefficient-field
cases passed $2114720$ complete assignments; twenty independent expanded
moments covered $6624$ assignments. Five larger acquisition-only cases
reach 77 binary residue coordinates, with no field scan during acquisition.
The raw 36-bit case reacquires the ring and phase, then the pair certificate
with $\gamma=10$, $\tau=0$, and multiplicities $(4,0)$. Its signed count
is again $3997696$.

That raw test obtains reference moments from an explicitly recorded
sixteen-element field scan and replays seven through the moment interface;
the zeroth is a Gaussian. A further raw 80-bit case at $d=4,e=4,b=5$
includes arbitrary higher cubic coefficients and rank-20 quotient noise.
It reacquires the pair template after its actual depth reduction and
agrees with an independently exhausted field leaf. An additional audit
checks 406 pair-Gaussian cases, covering 327056 phase evaluations and
every polar nullity $0,2,4$.

The production hypersurface adapter has also been checked by evaluating
112 actual equations over $\F_2$ and $\F_4$, including their
Artin--Schreier coordinates: $1728136$ complete assignments. The resulting
cardinalities agree with $432404$ separately expanded character evaluations.
The raw 36-bit callback test receives exactly the seven manufactured
coefficient dictionaries, with at most eight variables, and reproduces
the raw count. That callback transparently replays moments prepared by
the recorded sixteen-element field scan. For an opaque supplied backend,
the wrapper reports only its own enumeration count, not the backend's
unobserved internal work.

These executions validate acquisition, reductions and reconstruction.
They do not execute the generic Lauder--Wan engine. The raw wrapper
accepts a point counter receiving each manufactured hypersurface, or
labels its bounded reference scan. The earlier 49-moment construction
is preserved as an independently checked reference; the eight-moment
method is the production interface.

Quadratic-rank and finite-moment methods have established ancestry in
coding theory, including Luo \cite{luomoments}. For odd $d$, writing
$u=x+\omega y$ in $\F_{q^2}$ with $\omega^2+\omega+1=0$ identifies
these pair sums with a generalized Kasami sum and its linked linear
twist studied by Luo--Tang--Wang \cite{kasami}. The source proof gives
the exact identification. No invention of those classical sums or of
moment inversion is claimed. The candidate contribution is the
manufactured semilinear and block representation, the eight summaries
whose dimension is independent of unbounded module rank, and their
proved composition with raw weighted chain-ring counting. Historical
priority remains unestablished.
